\chapter{Painel Avançado II}
\label{cap:painel2}
% ── índice ──────────────────────────────────────────────────────
\index{Mundlak, estimador de|textbf}
\index{mundlak@\texttt{mundlak()}|textbf}
\index{efeitos correlacionados aleatórios|see{Mundlak, estimador de}}
\index{CRE|see{Mundlak, estimador de}}
\index{xtgls@\texttt{xtgls()}|textbf}
\index{modelo misto|textbf}
\index{mixedlm@\texttt{mixedlm()}|textbf}
\index{logit condicional|textbf}
\index{clogit@\texttt{clogit()}|textbf}
\index{REML}

% ─────────────────────────────────────────────────────────────────────────────
\section{Além do FE e RE básicos}

Os estimadores do cap.~\ref{cap:painel_avancado} (AB-GMM, PCSE, GEE) lidam
com dinâmica e heterocedasticidade. Aqui cobrimos quatro abordagens complementares:

\begin{itemize}
  \item \textbf{FE} — referência; efeitos fixos padrão (within estimator)
  \item \textbf{Mundlak (CRE)} — Correlated Random Effects; testa RE vs FE via $\gamma$
  \item \textbf{xtgls} — GLS para painel; flexível quanto à estrutura de $\Sigma$
  \item \textbf{mixedlm} — modelo misto com intercepts aleatórios
  \item \textbf{clogit} — logit condicional para variáveis binárias em painel
\end{itemize}

\subsection*{Estimador de Mundlak (CRE)}

Mundlak (1978) mostra que o estimador RE é equivalente ao FE se:
\[
  \alpha_i = \gamma'\,\bar{x}_i + u_i \quad u_i \perp x_{it}
\]
Regride $y_{it}$ em $x_{it}$ \emph{e} $\bar{x}_i$ (média por entidade).
Se $\hat{\gamma} \neq 0$, o efeito fixo é necessário.

\subsection*{GLS para painel (xtgls)}

O GLS de painel especifica uma estrutura para a covariância dos erros entre
entidades. Com \texttt{panels=hetero}, cada painel tem variância própria
$\sigma_i^2$, permitindo heterocedasticidade entre painéis.

\subsection*{Modelo misto (mixedlm)}

O modelo linear misto especifica:
\[
  y_{it} = \beta x_{it} + u_i + \varepsilon_{it} \quad u_i \sim N(0, \sigma_u^2)
\]
onde $u_i$ é um efeito aleatório de intercept por entidade. Difere do RE
padrão por usar REML e integrar sobre $u_i$ via máxima verossimilhança.

% ─────────────────────────────────────────────────────────────────────────────
\section{Verificação por DGP}

DGP: painel estático endógeno. $y = 1 + 2x + \alpha_i + \varepsilon$;
$\alpha_i = 0{,}05 \cdot i$ correlacionado com $x$ (via $z$).
$N = 50$ entidades, $T = 20$ períodos.

\begin{lstlisting}[caption={Painel II: FE, Mundlak, xtgls, mixedlm, clogit}]
seed(42)
generate df unit  = (_n - 1) % 50 + 1
generate df t     = (_n - (_n - 1) % 50 - 1) / 50 + 1
generate df alpha = unit * 0.05
generate df z     = (rnormal() - 0.5) * 1.7321
generate df x     = 0.5 * z + alpha + (rnormal() - 0.5)
generate df y     = 1.0 + 2.0 * x + alpha + e

let m_fe   = fe(y ~ x, df, id=unit)
let m_mund = mundlak(df, y ~ x, id="unit")
let m_gls  = xtgls(y ~ x, df, id=unit, time=t)
let m_mix  = mixedlm(y ~ x, df, id="unit")
generate df yb = (y > 2.5)
let m_cl   = clogit(yb ~ x, df, group="unit")
\end{lstlisting}

\begin{lstlisting}[language={}, basicstyle=\ttfamily\small, frame=none,
  numbers=none, backgroundcolor=\color{codbg}]
FE (within):        x = 2.041 ***    (verdadeiro: 2.0) ✓
Mundlak:            x = 2.040,  γ(x̄) = 0.939 ***  → rejeita RE
xtgls (hetero):     x = 2.768 ***    (tendencioso sem demeaning)
mixedlm:            x = 2.134 ***    (intercept aleatório: σ²_u = 0.388)
clogit (yb ~ x):    x = 1.397 ***    (grupos sem variação descartados)
\end{lstlisting}

\subsection*{Teste de Mundlak}

O teste $H_0: \gamma = 0$ (médias não correlacionadas com $x$) rejeita com
$F = 419{,}96$, $p < 0{,}001$. Isso confirma que $x$ é endógeno — o estimador
RE seria inconsistente e deve-se usar FE.

O xtgls sem remoção de FE superestima $\beta$ pois captura parte do efeito
individual. O mixedlm fica próximo do FE porque trata $\alpha_i$ como efeito
aleatório normal.

% ─────────────────────────────────────────────────────────────────────────────
\section{Sintaxe}

\begin{lstlisting}
fe(y ~ x, df, id=unit)                    // Fixed Effects (within)
mundlak(df, y ~ x, id="unit")             // CRE: inclui médias de grupo
xtgls(y ~ x, df, id=unit, time=t)         // GLS para painel
xtgls(y ~ x, df, id=unit, time=t, panels="corr")  // Σ correlacionada
mixedlm(y ~ x, df, id="unit")             // modelo misto (REML)
clogit(yb ~ x, df, group="unit")          // logit condicional
\end{lstlisting}

\begin{nota}
  No clogit, grupos onde $y$ não varia (todos 0 ou todos 1) são
  automaticamente descartados porque não contribuem para a verossimilhança
  condicional. Isso reduz a amostra efetiva, mas é teoricamente correto.
\end{nota}

\section{Painel não-linear via GEE}

\index{xtlogit@\texttt{xtlogit()}}
\index{xtprobit@\texttt{xtprobit()}}
\index{xtpoisson@\texttt{xtpoisson()}}
\index{GEE!painel não-linear}

Para modelos não-lineares em painel (logit, probit, Poisson), Hayashi
oferece wrappers sobre GEE (Generalized Estimating Equations), que
estimam modelos population-averaged com correlação intra-grupo
flexível:

\begin{lstlisting}
xtlogit(yb ~ x1 + x2, df, id="firm", corr="exchangeable")
xtprobit(yb ~ x1 + x2, df, id="firm", corr="exchangeable")
xtpoisson(y ~ x1 + x2, df, id="firm", corr="exchangeable")
\end{lstlisting}

Estruturas de correlação disponíveis:
\begin{itemize}
  \item \texttt{independence} — correlação intra-grupo zero
  \item \texttt{exchangeable} — simetria composta (padrão)
  \item \texttt{ar1} — autoregressivo de primeira ordem
  \item \texttt{unstructured} — matriz não-estruturada
\end{itemize}

\begin{nota}
  GEE estima efeitos populacionais (marginais), não efeitos condicionais
  aos efeitos fixos. Para logit com efeitos fixos, use \texttt{clogit}.
\end{nota}
